% ---------------------------------------------------------
% TABELA EMENTA 42: Empreendedorismo I
% ---------------------------------------------------------
\begin{xltabular}{\linewidth}{|X|p{2.5cm}|p{2.5cm}|}
\hline
\multirow{3}{=}{\textcolor{ifscgreen}{\textbf{Unidade Curricular:}}\\[3pt]\textbf{\large Empreendedorismo I}} & \multicolumn{2}{p{\dimexpr5cm+2\tabcolsep+\arrayrulewidth\relax}|}{\textcolor{ifscgreen}{\textbf{Semestre:}} \textbf{4º (Ano 2)}} \\ \cline{2-3}
 & \textcolor{ifscgreen}{\textbf{CH EaD*:}} & \textcolor{ifscgreen}{\textbf{CH Total*:}} \\ \cline{2-3}
 & \textbf{00 h} & \textbf{40 h} \\ \hline
\endhead
\multicolumn{3}{|p{\dimexpr\linewidth-2\tabcolsep-2\arrayrulewidth\relax}|}{\textcolor{ifscgreen}{\textbf{Objetivos:}}} \\ \hline
\multicolumn{3}{|p{\dimexpr\linewidth-2\tabcolsep-2\arrayrulewidth\relax}|}{
\begin{itemize}[leftmargin=*,noitemsep,topsep=2pt]
 \item \textit{[Objetivos de aprendizagem pendentes de preenchimento pelo docente responsável]}
\end{itemize}
} \\ \hline
\multicolumn{3}{|p{\dimexpr\linewidth-2\tabcolsep-2\arrayrulewidth\relax}|}{\textcolor{ifscgreen}{\textbf{Conteúdos:}}} \\ \hline
\multicolumn{3}{|p{\dimexpr\linewidth-2\tabcolsep-2\arrayrulewidth\relax}|}{
\begin{itemize}[leftmargin=*,noitemsep,topsep=2pt]
 \item \textit{[Conteúdo programático pendente de preenchimento pelo docente responsável]}
\end{itemize}
} \\ \hline
\multicolumn{3}{|p{\dimexpr\linewidth-2\tabcolsep-2\arrayrulewidth\relax}|}{\textcolor{ifscgreen}{\textbf{Estratégias de Ensino e Aprendizagem:}}} \\ \hline
\multicolumn{3}{|p{\dimexpr\linewidth-2\tabcolsep-2\arrayrulewidth\relax}|}{\textit{[Metodologia de ensino e avaliação pendente de preenchimento pelo docente responsável]}} \\ \hline
\multicolumn{3}{|p{\dimexpr\linewidth-2\tabcolsep-2\arrayrulewidth\relax}|}{\textcolor{ifscgreen}{\textbf{Bibliografia Básica:}}} \\ \hline
\multicolumn{3}{|p{\dimexpr\linewidth-2\tabcolsep-2\arrayrulewidth\relax}|}{1. \textit{[1º título da Bibliografia Básica pendente de indicação docente e validação no Parecer da Biblioteca do Câmpus]}
2. \textit{[2º título da Bibliografia Básica pendente de indicação docente e validação no Parecer da Biblioteca do Câmpus]}} \\ \hline
\multicolumn{3}{|p{\dimexpr\linewidth-2\tabcolsep-2\arrayrulewidth\relax}|}{\textcolor{ifscgreen}{\textbf{Bibliografia Complementar:}}} \\ \hline
\multicolumn{3}{|p{\dimexpr\linewidth-2\tabcolsep-2\arrayrulewidth\relax}|}{1. \textit{[1º título da Bibliografia Complementar pendente de indicação docente e validação no Parecer da Biblioteca do Câmpus]}
2. \textit{[2º título da Bibliografia Complementar pendente de indicação docente e validação no Parecer da Biblioteca do Câmpus]}
3. \textit{[3º título da Bibliografia Complementar pendente de indicação docente e validação no Parecer da Biblioteca do Câmpus]}} \\ \hline
\end{xltabular}

\vspace{0.4cm}


% ---------------------------------------------------------
% TABELA EMENTA 43: Responsabilidade Socioambiental e Sustentabilidade
% ---------------------------------------------------------
\begin{xltabular}{\linewidth}{|X|p{2.5cm}|p{2.5cm}|}
\hline
\multirow{3}{=}{\textcolor{ifscgreen}{\textbf{Unidade Curricular:}}\\[3pt]\textbf{\large Responsabilidade Socioambiental e Sustentabilidade}} & \multicolumn{2}{p{\dimexpr5cm+2\tabcolsep+\arrayrulewidth\relax}|}{\textcolor{ifscgreen}{\textbf{Semestre:}} \textbf{4º (Ano 2)}} \\ \cline{2-3}
 & \textcolor{ifscgreen}{\textbf{CH EaD*:}} & \textcolor{ifscgreen}{\textbf{CH Total*:}} \\ \cline{2-3}
 & \textbf{00 h} & \textbf{40 h} \\ \hline
\endhead
\multicolumn{3}{|p{\dimexpr\linewidth-2\tabcolsep-2\arrayrulewidth\relax}|}{\textcolor{ifscgreen}{\textbf{Objetivos:}}} \\ \hline
\multicolumn{3}{|p{\dimexpr\linewidth-2\tabcolsep-2\arrayrulewidth\relax}|}{
\begin{itemize}[leftmargin=*,noitemsep,topsep=2pt]
 \item Analisar os principais fatores responsáveis pela degradação socioambiental, considerando suas causas e consequências nos diferentes contextos.
 \item Compreender o conceito de sustentabilidade e suas aplicações.
 \item Aplicar os princípios da Responsabilidade Socioambiental e da Sustentabilidade na administração.
\end{itemize}
} \\ \hline
\multicolumn{3}{|p{\dimexpr\linewidth-2\tabcolsep-2\arrayrulewidth\relax}|}{\textcolor{ifscgreen}{\textbf{Conteúdos:}}} \\ \hline
\multicolumn{3}{|p{\dimexpr\linewidth-2\tabcolsep-2\arrayrulewidth\relax}|}{
\begin{itemize}[leftmargin=*,noitemsep,topsep=2pt]
 \item Relações entre sociedade, natureza e desenvolvimento.
 \item Principais impactos socioambientais das atividades humanas e das atividades produtivas.
 \item Desigualdade social e racismo ambiental.
 \item Atributos socioambientais do território e análise das potencialidades, vulnerabilidades e conflitos socioambientais.
 \item Histórico do conceito de desenvolvimento sustentável e sustentabilidade.
 \item Responsabilidade Socioambiental e programas de ESG (Ambiental, Social e Governança).
 \item Aplicações da sustentabilidade na administração.
\end{itemize}
} \\ \hline
\multicolumn{3}{|p{\dimexpr\linewidth-2\tabcolsep-2\arrayrulewidth\relax}|}{\textcolor{ifscgreen}{\textbf{Estratégias de Ensino e Aprendizagem:}}} \\ \hline
\multicolumn{3}{|p{\dimexpr\linewidth-2\tabcolsep-2\arrayrulewidth\relax}|}{- **Metodologia:** Desenvolvida em três blocos articulados: (1) levantamento das concepções prévias dos estudantes por meio de metodologias dialógicas, valorizando conhecimentos e vivências; (2) aprofundamento e contextualização dos temas a partir da realidade do território e das organizações, com pesquisas, estudos de caso, atividades práticas, trabalhos colaborativos e saídas de campo; (3) aplicação dos conhecimentos na elaboração de projetos voltados à área da Administração, contemplando responsabilidade socioambiental, sustentabilidade e gestão organizacional. São utilizados espaços como salas de aula, Sala Multidisciplinar, Laboratório de Meio Ambiente e Geomática, Laboratório de Informática e Centro Multiuso, com previsão de ao menos uma saída de campo. Não há necessidade de divisão de turma.
- **Avaliação:** Processual e contínua ao longo dos três blocos, com instrumentos como atividades em sala de aula ou entregues em datas estabelecidas, diário de bordo, pesquisas, trabalhos individuais e em grupo, apresentações orais, relatórios e desenvolvimento dos projetos.
- \textit{Nota: ementa editada e confirmada pela docente Elisa Gandolfo.}} \\ \hline
\multicolumn{3}{|p{\dimexpr\linewidth-2\tabcolsep-2\arrayrulewidth\relax}|}{\textcolor{ifscgreen}{\textbf{Bibliografia Básica:}}} \\ \hline
\multicolumn{3}{|p{\dimexpr\linewidth-2\tabcolsep-2\arrayrulewidth\relax}|}{1. CORREA, Fabio Marar Silveira. \textbf{Responsabilidade social empresarial e a sustentabilidade ambiental}. São Paulo: Dialética, 2023.
2. DIAS, Reinaldo. \textbf{Gestão ambiental: responsabilidade social e sustentabilidade}. 3. ed. São Paulo: Atlas, 2017.} \\ \hline
\multicolumn{3}{|p{\dimexpr\linewidth-2\tabcolsep-2\arrayrulewidth\relax}|}{\textcolor{ifscgreen}{\textbf{Bibliografia Complementar:}}} \\ \hline
\multicolumn{3}{|p{\dimexpr\linewidth-2\tabcolsep-2\arrayrulewidth\relax}|}{1. ALIGLERI, Lilian; ALIGLERI, Luiz Antonio; KRUGLIANSKAS, Isak. \textbf{Gestão socioambiental: responsabilidade e sustentabilidade do negócio}. São Paulo: Atlas, 2009.
2. COIMBRA, José de Ávila Aguiar (coord.). \textbf{Meio ambiente e consumismo}. São Paulo: Senac São Paulo, 2008. (Meio ambiente, 8).
3. LEITE, Paulo Roberto. \textbf{Logística reversa: sustentabilidade e competitividade}. 3. ed. São Paulo: Saraiva, 2017.
4. VEIGA, José Eli da. \textbf{Para entender o desenvolvimento sustentável}. São Paulo: Editora 34, 2015.} \\ \hline
\end{xltabular}

\vspace{0.4cm}


% ---------------------------------------------------------
% TABELA EMENTA 44: Empreendedorismo II
% ---------------------------------------------------------
\begin{xltabular}{\linewidth}{|X|p{2.5cm}|p{2.5cm}|}
\hline
\multirow{3}{=}{\textcolor{ifscgreen}{\textbf{Unidade Curricular:}}\\[3pt]\textbf{\large Empreendedorismo II}} & \multicolumn{2}{p{\dimexpr5cm+2\tabcolsep+\arrayrulewidth\relax}|}{\textcolor{ifscgreen}{\textbf{Semestre:}} \textbf{5º (Ano 3)}} \\ \cline{2-3}
 & \textcolor{ifscgreen}{\textbf{CH EaD*:}} & \textcolor{ifscgreen}{\textbf{CH Total*:}} \\ \cline{2-3}
 & \textbf{00 h} & \textbf{40 h} \\ \hline
\endhead
\multicolumn{3}{|p{\dimexpr\linewidth-2\tabcolsep-2\arrayrulewidth\relax}|}{\textcolor{ifscgreen}{\textbf{Objetivos:}}} \\ \hline
\multicolumn{3}{|p{\dimexpr\linewidth-2\tabcolsep-2\arrayrulewidth\relax}|}{
\begin{itemize}[leftmargin=*,noitemsep,topsep=2pt]
 \item Compreender o processo empreendedor e as etapas do planejamento de novos empreendimentos.
 \item Identificar e avaliar oportunidades de negócio.
 \item Desenvolver um modelo de negócios, aplicando técnicas de ideação, modelagem, validação e prototipação.
 \item Aplicar técnicas de comunicação de ideias inovadoras.
\end{itemize}
} \\ \hline
\multicolumn{3}{|p{\dimexpr\linewidth-2\tabcolsep-2\arrayrulewidth\relax}|}{\textcolor{ifscgreen}{\textbf{Conteúdos:}}} \\ \hline
\multicolumn{3}{|p{\dimexpr\linewidth-2\tabcolsep-2\arrayrulewidth\relax}|}{
\begin{itemize}[leftmargin=*,noitemsep,topsep=2pt]
 \item Identificação e avaliação de oportunidades de negócios.
 \item Modelagem, validação e desenvolvimento de novos negócios.
 \item Estruturação de modelos e planos de negócios.
 \item Pitch de negócios: estrutura narrativa e técnicas de apresentação.
 \item Formalização de empreendimentos e ecossistema de apoio ao empreendedorismo e inovação.
\end{itemize}
} \\ \hline
\multicolumn{3}{|p{\dimexpr\linewidth-2\tabcolsep-2\arrayrulewidth\relax}|}{\textcolor{ifscgreen}{\textbf{Estratégias de Ensino e Aprendizagem:}}} \\ \hline
\multicolumn{3}{|p{\dimexpr\linewidth-2\tabcolsep-2\arrayrulewidth\relax}|}{- **Metodologia:** Fundamentada em metodologias ativas, com ênfase na aprendizagem baseada em projetos, oficinas práticas, desenvolvimento colaborativo de modelos de negócios, prototipação, validação de ideias e apresentação de propostas de empreendimentos, preferencialmente relacionadas às demandas do território e aos desafios sociais, ambientais e econômicos da comunidade. Utiliza sala de aula, laboratório de informática, sala multidisciplinar e centro multiuso. Articula-se com Empreendedorismo I, Projeto Integrador I e II, Responsabilidade Socioambiental e Sustentabilidade, línguas inglesa, portuguesa e espanhola, Artes, Sociologia e Filosofia.
- **Avaliação:** Diagnóstica, processual e formativa, com instrumentos diversificados como atividades individuais e em grupo, estudos de caso, resolução de problemas, pesquisas, elaboração de canvas, relatórios, apresentações, portfólios, protótipos, modelos de negócios, pitches, autoavaliação, avaliação por pares e avaliações escritas, quando pertinentes.} \\ \hline
\multicolumn{3}{|p{\dimexpr\linewidth-2\tabcolsep-2\arrayrulewidth\relax}|}{\textcolor{ifscgreen}{\textbf{Bibliografia Básica:}}} \\ \hline
\multicolumn{3}{|p{\dimexpr\linewidth-2\tabcolsep-2\arrayrulewidth\relax}|}{1. BESSANT, John; TIDD, Joe; COSTA, Francisco Araújo da. \textbf{Inovação e empreendedorismo}. 3. ed. Porto Alegre: Bookman, 2019.
2. DORNELAS, José Carlos Assis. \textbf{Empreendedorismo corporativo: como ser empreendedor, inovar e se diferenciar na sua empresa}. 2. ed. Rio de Janeiro: Elsevier, 2008.
3. DORNELAS, José Carlos Assis. \textbf{Empreendedorismo: transformando idéias em negócios}. 3. ed. rev. e atual. Rio de Janeiro: Elsevier, 2008.
4. HISRICH, Robert D.; PETERS, Michael P.; SHEPHERD, Dean A. \textbf{Empreendedorismo}. Tradução de Francisco Araujo da Costa. 9. ed. Porto Alegre: AMGH, 2014.
5. OLIVEIRA, Edson Marques. \textbf{Empreendedorismo social: da teoria à prática, do sonho à realidade}. Rio de Janeiro: Qualitymark, 2008.} \\ \hline
\multicolumn{3}{|p{\dimexpr\linewidth-2\tabcolsep-2\arrayrulewidth\relax}|}{\textcolor{ifscgreen}{\textbf{Bibliografia Complementar:}}} \\ \hline
\multicolumn{3}{|p{\dimexpr\linewidth-2\tabcolsep-2\arrayrulewidth\relax}|}{1. BERNARDI, Luiz Antonio. \textbf{Manual de empreendedorismo e gestão: fundamentos, estratégias e dinâmicas}. 2. ed. São Paulo: Atlas, 2012.
2. DORNELAS, José Carlos Assis. \textbf{Empreendedorismo corporativo: como ser empreendedor, inovar e se diferenciar na sua empresa}. 2. ed. Rio de Janeiro: Elsevier, 2008.
3. OLIVEIRA, Djalma de Pinho Rebouças de. \textbf{Empreendedorismo: vocação, capacitação e atuação direcionadas para o plano de negócios}. São Paulo: Atlas, 2014.
4. OSTERWALDER, Alexander; PIGNEUR, Yves. \textbf{Business model generation = inovação em modelos de negócios: um manual para visionários, inovadores e revolucionários}. Rio de Janeiro: Alta Books, 2011.} \\ \hline
\end{xltabular}

\vspace{0.4cm}

